\documentclass[11pt,letterpaper]{article}
\usepackage[margin=0.83in]{geometry}
\usepackage{amsmath,amssymb,amsthm,mathtools,microtype}
\usepackage[hidelinks]{hyperref}
\usepackage{enumitem}
\setlist{nosep,leftmargin=*}
\newtheorem{theorem}{Theorem}
\newtheorem{corollary}[theorem]{Corollary}
\newtheorem{proposition}[theorem]{Proposition}
\newcommand{\HH}{\mathcal H}
\newcommand{\BB}{\mathcal B}
\newcommand{\CP}{\mathbb C_+}
\newcommand{\Rea}{\operatorname{Re}}
\newcommand{\ord}{\operatorname{ord}}
\newcommand{\dd}{\mathrm d}
\title{Twistor bridge continuation}
\author{Research continuation for Jude Gomila}
\date{3 October 2026}
\begin{document}
\maketitle
\vspace{-1em}
\noindent\textbf{Results.} The exact wide split-twistor bridge cannot be divided into identical holomorphic operator stages, even with angular mixing. The actual completed arithmetic ratio has a negative delay kernel before the first prime enters, ruling out a positive scalar convolution model for every thin width. A constructive continuation nevertheless gives compact arithmetic Hankel operators at every positive width, and an explicit nontrivial window with norm at most $2^{-a}<1$. A further interval certificate proves norm at most $0.99998$ at width $a=1/4$ for every truncation $L\leq3/2$, crossing the first prime threshold. Global spectral control remains unresolved. None of these results proves RH or excludes other twistor constructions.

\section{Established construction and the actual gap}
Write
\[
 \xi(s)=\tfrac12s(s-1)\pi^{-s/2}\Gamma(s/2)\zeta(s),\qquad
 X(w)=\xi(\tfrac12+w),\qquad R_a(w)=\frac{X(w-a)}{X(w+a)}.
\]
The split-twistor construction in the 1 October report [1] acts on
\[
 \HH=L^2_{\rm even}(S^3)=\widehat{\bigoplus}_{m\geq0}\HH_m,
 \qquad \HH_m=\mathcal Y_{2m}(S^3),\quad d_m=\dim\HH_m=(2m+1)^2.
\]
Its normalized Fourier part and completed arithmetic bridge are
\begin{equation}\label{eq:model}
 A(w)|_{\HH_m}=b_m(w)I,\quad
 b_m(w)=\prod_{j=1}^{m}\frac{w-2j}{w+2j},\qquad
 T(w)=R_{3/2}(w)A(w).
\end{equation}
The empty product is one. Both are holomorphic contractions on $\CP=\{\Rea w>0\}$. The arithmetic identity is
\[
 R_{3/2}(w)=\frac{\xi(w-1)}{\xi(w+2)}
 =\frac{(w-1)(w-2)}{(w+1)(w+2)}
 \frac{L(2-w)}{L(2+w)},\qquad L(s)=\pi^{-s/2}\Gamma(s/2)\zeta(s).
\]
It is derived by real and prime-place Fourier factors, initially in $|\Rea w|<1$, followed by meromorphic continuation. The geometry is $\mathbb{RP}^3$ and its dual, with $SL(4,\mathbb R)$ acting on homogeneous degree $-2-w$ distributions. The parameter reflection is $w\mapsto-w$; the fixed thickness $3/2$ is not the homogeneity parameter.

The wide-shift contractivity is unconditional [2]. For the candidate $T_a=R_aA$, the constant harmonic is an eigenvector with eigenvalue $R_a$. Hence $\|T_a(w)\|=|R_a(w)|$ at regular points in $\CP$. Proving that these families are Schur for every $a>0$, or a sequence $a\downarrow0$, remains equivalent to RH [1,3]. The present investigation asks whether a passive wide bridge can be subdivided to manufacture that missing family.

\section{A local obstruction without a symmetry assumption}
At the positive real point $w=2$, exact completion gives
\begin{equation}\label{eq:degeneration}
 R_{3/2}(2)=\frac{15}{2\pi^2}>0,\qquad
 T(2)=\frac{15}{2\pi^2}P_0,\qquad A(2)=P_0,
\end{equation}
where $P_0$ is the projection onto the constants. Each nonconstant multiplier has a simple zero there:
\begin{equation}\label{eq:slope}
 b'_m(2)=\frac{(-1)^{m-1}}{2m(m+1)}\ne0\quad(m\geq1).
\end{equation}
Indeed, differentiate the first factor, whose derivative at 2 is $1/4$, and telescope the remaining product $\prod_{j=2}^{m}(1-j)/(1+j)$.

\begin{theorem}[No holomorphic operator root]\label{thm:root}
Let $B=A$ or $B=T$. For every integer $q\geq2$, there is no neighborhood $D$ of $2$ and holomorphic map $Q:D\to\BB(\HH)$ such that
\[
 Q(w)^q=B(w)\qquad(w\in D).
\]
No rotation equivariance, normality, or contractivity of $Q$ is assumed.
\end{theorem}
\begin{proof}
Shrink $D$ to a disk of radius less than 1 lying in $\CP$, with $R_{3/2}$ nonzero on $D$. Set $r=1$ for $B=A$, or $r=R_{3/2}$ for $B=T$, and write $t_m=rb_m$. For $w\in D\setminus\{2\}$ all $t_m$ are nonzero and
\[
 |b_{m+1}(w)|=|b_m(w)|\left|\frac{w-2m-2}{w+2m+2}\right|<|b_m(w)|.
\]
The strict inequality follows from
$|w+c|^2-|w-c|^2=4c\Rea w>0$ for $c>0$. Thus the eigenvalues $t_m(w)$ are pairwise distinct, and their eigenspaces are exactly $\HH_m$.

If $Q^q=B$, then $QB=BQ$. For $v\in\HH_m$, this gives $B(Qv)=t_mQv$, so $Qv\in\HH_m$ on the punctured disk. Continuity extends this invariance to $w=2$. Therefore $Q_m=Q|_{\HH_m}$ is a holomorphic $d_m\times d_m$ matrix and
\[
 Q_m(w)^q=t_m(w)I_{d_m},\qquad (\det Q_m(w))^q=t_m(w)^{d_m}.
\]
For $m\geq1$, $t_m$ has a simple zero at 2, by \eqref{eq:degeneration}--\eqref{eq:slope}. Taking zero orders yields
\[
 q\,\ord_{w=2}(\det Q_m)=d_m.
\]
In particular $q$ divides both $d_1=9$ and $d_2=25$. Since their greatest common divisor is 1, this is impossible for $q\geq2$.
\end{proof}

The proof is local, uses no information about nontrivial zeta zeros, and remains valid for $r(w)A(w)$ whenever $r$ is holomorphic and nonzero at 2. Therefore changing to any nonsingular scalar normalization there does not fix the obstruction. It also survives an invertible holomorphic similarity transformation. Even a globally continuous operator root on a punctured disk is impossible: taking determinant winding numbers around the puncture gives the same divisibility contradiction.

Deleting the constant harmonic does not help. Nor does deleting finitely many low angular modes: any two remaining adjacent nonconstant harmonic spaces have coprime dimensions $(2m+1)^2$ and $(2m+3)^2$, and the same argument applies. The obstruction is therefore present in the geometric sector itself, independently of the arithmetic constant mode.

\section{Finite enlargement gives a sharp boundary}
\begin{theorem}[Exact amplification criterion]\label{thm:amp}
For a positive integer $r$, the amplified bridge $B(w)\otimes I_r$ admits a holomorphic $q$th root near $2$ if and only if $q$ divides $r$. When $q\mid r$, a root holomorphic and contractive on all of $\CP$ exists.
\end{theorem}
\begin{proof}
For necessity, the preceding proof has eigenspace dimensions $9r$ and $25r$, so $q$ divides their greatest common divisor $r$.

For sufficiency, first take $r=q$ and define the cyclic block operator on $\HH^{\oplus q}$ by
\begin{equation}\label{eq:companion}
 C_q(B)(x_1,\ldots,x_q)=(Bx_q,x_1,\ldots,x_{q-1}).
\end{equation}
After $q$ steps each coordinate has passed through $B$ exactly once, so $C_q(B)^q=B^{\oplus q}$. It is holomorphic and
\[
 \|C_q(B)x\|^2=\|Bx_q\|^2+\sum_{j=1}^{q-1}\|x_j\|^2\leq\|x\|^2.
\]
For $r=qk$, take $k$ copies of this construction.
\end{proof}

\begin{corollary}[No finitely amplified thickness semigroup]
For fixed finite $r$, there is no family $S_t(w)$, $t\geq0$, of holomorphic bounded operators on $\HH\otimes\mathbb C^r$ satisfying
\[
 S_{t+u}(w)=S_t(w)S_u(w),\qquad S_1(w)=B(w)\otimes I_r.
\]
In particular a continuous passive semigroup with this exact endpoint cannot exist on that space.
\end{corollary}
\begin{proof}
Choose an integer $q\geq2$ that does not divide $r$. Then $S_{1/q}^q=S_1$, contradicting Theorem~\ref{thm:amp}. No continuity in $t$ was needed.
\end{proof}

\textbf{Scope.} This excludes identical-stage pointwise subdivision of this transfer operator, including arbitrary finite copies of its complete angular space. It does not exclude every finite-dimensional auxiliary construction with a different endpoint, a compression, different source and target spaces, unbounded maps, or a family without the semigroup law. Unequal passive cascades are also not ruled out.

\section{Infinite memory permits subdivision but supplies no arithmetic}
The obstruction is not an assertion that passive semigroups cannot be built. An explicit construction shows exactly what an enlargement can do. On $\mathcal K=L^2([0,1);\HH)$, for $t\geq0$ and $0\leq x<1$, set
\begin{equation}\label{eq:delay}
 (U_t(w)f)(x)=B(w)^{n(x,t)}f\bigl(x+t-n(x,t)\bigr),
 \qquad n(x,t)=\lfloor x+t\rfloor.
\end{equation}
Values at the finitely many partition boundaries are immaterial in $L^2$.

\begin{proposition}[Explicit passive delay enlargement]
The operators in \eqref{eq:delay} form a strongly continuous contraction semigroup in $t$. For each fixed $t$, they are holomorphic in $w\in\CP$, and $U_1(w)=I_{L^2([0,1))}\otimes B(w)$.
\end{proposition}
\begin{proof}
Adding $t$ and $u$ adds the number of wraps around the interval, so powers of $B$ multiply and $U_tU_u=U_{t+u}$. The map $x\mapsto x+t\pmod1$ preserves Lebesgue measure, and $\|B^n\|\leq1$; hence $\|U_tf\|\leq\|f\|$. For fixed $t$, $n(x,t)$ takes at most two consecutive integer values, and its partition is independent of $w$. Thus $U_t(w)$ is a finite sum of fixed bounded translation and restriction operators multiplied by powers of the holomorphic family $B(w)$.

Strong continuity at $t=0$ follows from continuity of translations in $L^2$, while the wrapped interval has measure $t$ and its contribution tends to zero for each $f$. The semigroup law then gives strong continuity at every $t\geq0$. At $t=1$ every point wraps once, proving the endpoint identity.
\end{proof}

This infinite-dimensional auxiliary coordinate evades the finite-multiplicity determinant argument. It applies to \emph{every} holomorphic contraction $B$, so its existence contains no special fact about the Riemann arithmetic. Neither $C_q(B)$ nor $U_t$ has been identified with $T_{(3/2)t}$; merely assigning it that name would not establish the identity.

\section{The exact condition a successful continuation must add}
The shifted scalar ratios do not obey the pointwise semigroup law. They instead satisfy the meromorphic cocycle
\begin{equation}\label{eq:cocycle}
 R_{a+b}(w)=R_a(w-b)\,R_b(w+a).
\end{equation}
This is immediate by cancellation of the middle factors $X(w+a-b)$. However $w-b$ can lie outside $\CP$ even when $w\in\CP$, so a Schur bound on each right half-plane does not transfer through this identity to a proof of thinning. Extracting ordinary roots is not the correct algebraic operation for these arithmetic shifts.

A constructive continuation therefore needs more than passive geometric subdivision. It must derive a bounded input/output coefficient equal to
\[
 R_a(w)=\frac{\xi(1/2+w-a)}{\xi(1/2+w+a)}
\]
for arbitrarily small positive $a$, from a positive realization justified independently of that coefficient's Schur property. If the coefficient is recovered by unit-norm input/output vectors, passivity would give $|R_a(w)|\leq1$. Holomorphicity must also be proved, rather than inserted by assuming the absence of its possible poles.

Equivalently, the undivided arithmetic target remains
\[
 D(x,t)=\Rea\!\left(X'(x+it)\overline{X(x+it)}\right)
       =\tfrac12\partial_x|X(x+it)|^2\geq0
 \quad(x>0,\ t\in\mathbb R).
\]
The conjugate in this formula is essential. A right-half-plane zero of multiplicity $m$ would make $D$ negative immediately to its left, so the target is not a harmless choice of normalization.

\textbf{Current conclusion.} The split twistor and dual-twistor transform, its prime completion, and its wide operator are explicit. The added result is a proved obstruction to identical-stage geometric thinning, together with sharp constructions showing its limits. The arithmetic positivity theorem is still open. This is a useful restriction on the next construction, not a new zero-free region.

\section{The actual arithmetic ratio has a negative delay interval}
The scalar arithmetic obstruction can be made more specific than a generic counterexample. Use a separate complex variable $z$, with
\[
 F_a(z)=\frac{\xi(z-a)}{\xi(z+a)}=R_a(z-\tfrac12),\qquad
 \phi_{a,\sigma}(t)=\frac{F_a(\sigma+it)}{F_a(\sigma)},\qquad \sigma>1+a.
\]
The variables $a$ and $\sigma$ in this section are width and Euler-domain real part; neither is the angular spectral variable $w$ at the root obstruction.

\begin{theorem}[No positive scalar delay model at thin widths]\label{thm:negative}
For $0<a\leq1$ and $\sigma>1+a$, $\phi_{a,\sigma}$ is not the characteristic function of any probability measure. In particular, this excludes a positive scalar convolution-kernel realization even in the absolutely convergent Euler domain. It does not exclude Schur contractivity or a positive Hilbert-space energy realization.
\end{theorem}
\begin{proof}
The Euler expansion is
\[
 \frac{\zeta(z-a)}{\zeta(z+a)}=\sum_{n\geq1}c_a(n)n^{-z},\qquad
 c_a(n)=n^a\prod_{p\mid n}(1-p^{-2a})\geq0,\quad c_a(1)=1.
\]
Define, for $0<y<1$,
\begin{equation}\label{eq:g}
 g_a(y)=\frac{2\pi^a}{\Gamma(a)}\left[
 y^{2-a}(1-y^2)^{a-1}
 -a y^{a-1}\int_{y^2}^{1}t^{1/2-a}(1-t)^{a-1}\,\dd t\right].
\end{equation}
This is Suzuki's archimedean kernel [3, Section 2]. A direct beta integral and interchange of the two integrals verify
\[
 \int_0^1 g_a(y)y^{z-1}\,\dd y
 =\pi^a\frac{(z-a)(z-a-1)}{(z+a)(z+a-1)}
   \frac{\Gamma((z-a)/2)}{\Gamma((z+a)/2)}.
\]
Consequently, the full additive Laplace kernel is
\begin{equation}\label{eq:q}
 q_a(u)=\sum_{n\leq e^u}c_a(n)g_a(ne^{-u}),\qquad
 F_a(z)=\int_0^\infty e^{-zu}q_a(u)\,\dd u
 \quad(\Rea z>1+a).
\end{equation}
Threshold values, which may be singular, are interpreted almost everywhere. There is no extra factor $e^{-u}$ in $q_a$.

To prove a negative interval, put $M=(1-y^2)/y^2$. The bracket in \eqref{eq:g}, divided by $y^{a-1}$, is
\[
 yM^{a-1}-a\int_0^M v^{a-1}(1+v)^{-3/2}\,\dd v.
\]
If $U_a$ has density $au^{a-1}$ on $(0,1)$, then $U_a$ is stochastically smaller than a uniform variable for $0<a\leq1$. Since $(1+Mu)^{-3/2}$ is decreasing,
\[
 \frac{a}{M^a}\int_0^M v^{a-1}(1+v)^{-3/2}\,\dd v
 \geq\int_0^1(1+Mv)^{-3/2}\,\dd v
 =\frac{2(1-y)}{M}.
\]
The bracket in \eqref{eq:g} is therefore at most $y^{a-1}M^{a-1}(3y-2)$, and $g_a(y)<0$ for $y<2/3$. On the prime-free interval
\begin{equation}\label{eq:negativeinterval}
 \log(3/2)<u<\log2,
\end{equation}
only $n=1$ occurs in \eqref{eq:q}, so $q_a(u)=g_a(e^{-u})<0$.

The measure $e^{-\sigma u}q_a(u)\dd u$ has finite total variation. For completeness, beta inversion writes $g_a(e^{-u})$ as a positive integrable beta kernel minus its convolution with $2ae^{-(a-1)u}$. Its damped absolute integral is bounded by
\[
 \pi^a\frac{\Gamma((\sigma-a)/2+1)}{\Gamma((\sigma+a)/2+1)}
 \left(1+\frac{2a}{\sigma+a-1}\right).
\]
Multiply this bound by the convergent positive series $\sum c_a(n)n^{-\sigma}$ to obtain finite total variation for the full kernel. Since $F_a(\sigma)>0$, the normalized signed measure
\[
 F_a(\sigma)^{-1}e^{-\sigma u}q_a(u)\,\dd u
\]
has mass one and Fourier transform $\phi_{a,\sigma}$, using the $e^{-itu}$ convention. It has negative mass on \eqref{eq:negativeinterval}. Uniqueness of the Fourier transform of finite measures rules out any alternative positive probability measure.
\end{proof}

At $a=1$ the inversion is elementary:
\[
 F_1(z)=2\pi\frac{z-2}{z(z+1)}\frac{\zeta(z-1)}{\zeta(z+1)},\qquad
 g_1(e^{-u})=2\pi(3e^{-u}-2).
\]
Changing Euler or Dirichlet coefficients with $n\geq2$ cannot change the negative interval, since those delays begin at $\log2$. A successful construction must allow a signed impulse response or a different, genuinely coupled state realization. Positive arithmetic coefficients alone are insufficient after the exact $z(z-1)$ completion.

\section{Compact arithmetic operators and an explicit local spectral gap}
The negative scalar kernel does not prevent positive operator energy forms. Define Suzuki's multiplicative kernel and truncated Hankel operator by
\[
 h_a(x)=x^{-1}q_a(\log x)\ (x>1),\qquad h_a(x)=0\ (x<1),\qquad
 (\mathsf H_{a,L}f)(x)=\int_0^Lh_a(xy)f(y)\,\dd y.
\]
Here $a$ is the arithmetic width and $L$ is the truncation endpoint. The endpoint is distinct from the width parameter used in \eqref{eq:cocycle}.

\begin{proposition}[Compactness at every positive width]\label{prop:compact}
For every fixed $a>0$ and $L>1$, $\mathsf H_{a,L}$ is a bounded compact self-adjoint operator on $L^2((0,L),\dd x)$. It is not Hilbert--Schmidt when $0<a\leq1/2$.
\end{proposition}
\begin{proof}
The support condition $xy\geq1$ confines the nonzero part to $[1/L,L]^2$. Only finitely many thresholds $xy=n\leq L^2$ occur. The beta formula gives a locally integrable singularity of order $(xy-n)^{a-1}$ on its supported side, plus bounded or smoother terms. If $0<a<1$, remove strips $|xy-n|<\varepsilon$ around these thresholds. The remaining bounded kernel defines a Hilbert--Schmidt operator. Since $x,y$ stay bounded away from zero, the absolute row and column integrals of the removed singular parts are $O(\varepsilon^a)$; the Schur test proves norm convergence as $\varepsilon\downarrow0$. For $a\geq1$ the kernel is bounded on this compact square and the claim is immediate. Reality and symmetry in $x,y$ give self-adjointness.

At $xy=1$, the nonzero leading coefficient of the first singularity gives square magnitude proportional to $(xy-1)^{2a-2}$. Its two-dimensional integral diverges for $a\leq1/2$, so the kernel is not square integrable and the operator is not Hilbert--Schmidt.
\end{proof}

Let $\gamma$ denote Euler's constant, and set
\begin{equation}\label{eq:window}
 \ell_0=\frac{e^{-\gamma}}{8\pi},\qquad L_0=e^{\ell_0}.
\end{equation}
\begin{theorem}[A common nontrivial window for all thin widths]\label{thm:local}
For every $0<a\leq1$,
\[
 \|\mathsf H_{a,L_0}\|\leq2^{-a}<1,\qquad
 I\pm\mathsf H_{a,L_0}\succeq(1-2^{-a})I.
\]
In particular both signs are invertible, with inverse norm at most $(1-2^{-a})^{-1}$. No RH assumption is used.
\end{theorem}
\begin{proof}
On its nonzero subspace, the unitary map $(\mathcal Uf)(v)=e^{v/2}f(e^v)$ takes the operator to the interval $[-\ell_0,\ell_0]$ with kernel
\[
 K_a(v+v')=e^{-(v+v')/2}g_a(e^{-(v+v')})\,\mathbf1_{v+v'>0}.
\]
There is only the $n=1$ term, since $e^{2\ell_0}<2$.
The kernel is nonnegative on this interval. Indeed, for $y=e^{-u}$ and $M=(1-y^2)/y^2$,
$a\int_0^Mv^{a-1}(1+v)^{-3/2}\dd v\leq M^a$, so the bracket in \eqref{eq:g} is at least $y^{a-1}M^{a-1}(y-M)$. Here $0<u\leq2\ell_0<1/12$, so $y>1-u>11/12$ and $y^3+y^2>1$, hence $y>M$.

For $0<a\leq1$, discard the subtractive term of \eqref{eq:g} and use $1-e^{-2u}\geq2ue^{-2u}$ to obtain
\[
 0\leq K_a(u)\leq g_a(e^{-u})
 \leq\frac{(2\pi)^a}{\Gamma(a)}u^{a-1}\quad(0<u\leq2\ell_0).
\]
Every row and column integral is therefore at most
\[
 \frac{(4\pi\ell_0)^a}{\Gamma(1+a)}
 \leq(4\pi e^\gamma\ell_0)^a=2^{-a}.
\]
The second inequality follows from convexity of $\log\Gamma$, namely $\log\Gamma(1+a)\geq-\gamma a$. The Schur test gives the norm bound, and self-adjointness gives the two positive energy forms and inverse estimates.
\end{proof}

\textbf{Positivity distinction.} A nonnegative kernel preserves nonnegative functions but need not define a positive quadratic form. Here the Hankel operator is indefinite: a bump at negative logarithmic coordinates has zero self-pairing but can have positive cross-pairing with a bump at positive coordinates. Subtracting a sufficiently small multiple of the latter gives a negative quadratic form. The positive operators proved above are $I+\mathsf H$ and $I-\mathsf H$.

There is a structural reason the inverse estimate cannot be uniform at zero width. As $a\downarrow0$, the positive kernels $K_a$ have mass tending to 1 on $(0,2\ell_0)$, and mass tending to zero outside every neighborhood of 0. The beta-kernel formula proves this; the subtractive convolution has mass $O(a)$ on a fixed compact interval. Thus they form a one-sided approximate identity, and the logarithmic operators converge strongly to reflection $Jf(v)=f(-v)$. Their norms tend to 1. Even and odd test functions show that neither inverse family $I\pm\mathsf H$ can remain uniformly bounded. Convergence is not in operator norm, since each truncated operator is compact while $J$ is not.

\textbf{What this does and does not close.} Compactness removes a local operator-definition obstacle below the Hilbert--Schmidt threshold. The explicit norm bound proves both Fredholm inverses on a common nonzero window. It does not supply the canonical Hamiltonian or its global positivity; ordinary Fredholm determinant formulas require additional justification below their trace-ideal thresholds. The bound tends to 1 as $a\downarrow0$, and the inverse estimate grows like $1/(a\log2)$.

The first nontrivial prime threshold is $L=\sqrt2$, where $xy=2$ enters. Sign changes already occur in the $n=1$ kernel whenever entries satisfy $3/2<xy<2$, so a nonnegative-kernel argument cannot cover all larger windows. Neither fact by itself disproves $\|\mathsf H_{a,L}\|<1$; the next missing theorem is spectral control for the signed, coupled arithmetic kernel as $L$ increases, with quantified behavior as $a\downarrow0$. The local propositions are derived here from the explicit classical kernel; their literature novelty has not been assessed.

\section{A certified obstruction after the first prime threshold}
The next bounded test used $a=1/4$ and $L=3/2$, beyond the first prime threshold $L=\sqrt2$. Put $V=2\log(3/2)$ and $b=\log2$. In logarithmic coordinates the signed kernel on $[-V/2,V/2]$ is
\[
 K(v)=e^{-v/2}\left[g_{1/4}(e^{-v})
 +(2^{1/4}-2^{-1/4})g_{1/4}(2e^{-v})\mathbf1_{v>b}\right]\mathbf1_{v>0}.
\]
Only $n=1,2$ occur. Let $\mathsf H_{|K|}$ denote the operator obtained by replacing $K$ with $|K|$. For the normalized constant test function, and any measurable $\eta$ with $|\eta|\leq1$,
\begin{equation}\label{eq:witness}
 \|\mathsf H_{|K|}\|\geq\frac1V\int_0^V(V-v)|K(v)|\,\dd v
 \geq\frac1V\int_0^V(V-v)\eta(v)K(v)\,\dd v.
\end{equation}
Choose the successive signs $+1,-1,+1,-1$ on the intervals with cuts
\[
 0,\quad 3/10,\quad\log2,\quad3/4,\quad V.
\]
No root isolation or assumption about the sign of $K$ is needed for the second inequality.

\begin{proposition}[Computer-assisted absolute-kernel obstruction]
For this exact width and window, the rightmost integral in \eqref{eq:witness} lies in
\[
 [1.0494529987629657540,\ 1.0651499622678697025].
\]
Consequently $\|\mathsf H_{|K|}\|>1$. Methods that discard the kernel's signs and then bound the absolute-kernel operator by 1 cannot establish contractivity here.
\end{proposition}
\begin{proof}[Certificate]
The script \texttt{certify\_absolute\_obstruction.py} uses outward-rounded interval arithmetic at 40 decimal digits. It splits at the displayed cuts and the prime threshold. For each component starting at its singularity, the substitution $v=b_j+d x^4$ cancels the endpoint power $t^{-3/4}$ exactly. The resulting continuous integrands are enclosed on 128 rectangles per segment and component.

The incomplete beta term is evaluated with
\[
 B_z(1/4,5/4)=z^{1/4}\sum_{k\geq0}\frac{c_kz^k}{k+1/4},\qquad
 c_0=1,\quad c_k=c_{k-1}\frac{k-1-1/4}{k}.
\]
All $c_k$ for $k\geq1$ are negative. Here $0\leq z\leq65/81$ exactly. The omitted series after degree 40 has absolute value at most
\[
 \frac{|c_{41}|}{41+1/4}\frac{(65/81)^{41}}{1-65/81}
 <2.924\times10^{-8}.
\]
The script encloses that remainder with its correct negative sign. It uses the decreasing function $(1-e^{-2t})/t$, with its limit 2 at zero, to avoid an undefined endpoint division. Summing the interval integrals gives the displayed enclosure. The complete executable certificate and interval output are supplied in the reproducible bundle.
\end{proof}

\textbf{Numerical motivation for the signed certificate.} Singularity-integrated, 1024-cell Galerkin compressions at width $1/4$ gave the numerical norms
\[
 \begin{array}{c|ccc}
 L&1.3&\sqrt2&1.5\\\hline
 \|P_N\mathsf H P_N\|\text{ (diagnostic)}
 &0.993768391046&0.999654874214&0.999961720777.
 \end{array}
\]
These are finite-compression diagnostics, not validated upper bounds on the full signed operator. The elementary singular-strip remainder bound is much larger than the apparent gaps.

A direct block certificate would retain cancellation as follows. For a projection $P$ and $Q=I-P$, validated bounds
\[
 \|P\mathsf H P\|\leq\alpha,\quad
 \|Q\mathsf H Q\|\leq\beta<1,\quad
 \|Q\mathsf H P\|\leq\varepsilon,\qquad
 \alpha+\frac{\varepsilon^2}{1-\beta}<1
\]
would imply positivity of both $I\pm\mathsf H$ by the block Schur complement. This requires the complete cross-block norm, not just convergence of Ritz values or the residual of one vector. The elementary cell-projection estimate did not close this condition. The next section supplies a sharper Fourier-dilation Schur certificate for the specific window $a=1/4$, $L=3/2$.

\section{A validated signed certificate beyond the first prime}
\begin{theorem}[Computer-assisted signed contraction]\label{thm:signed}
At the fixed width $a=1/4$, for every $1<L\leq3/2$,
\[
 \|\mathsf H_{1/4,L}\|\leq\frac{49999}{50000}=0.99998,
 \qquad I\pm\mathsf H_{1/4,L}\succeq\frac1{50000}I.
\]
This is a full infinite-dimensional operator bound. It uses no hypothesis about zeta zeros. The arithmetic and matrix enclosures use outward-rounded \texttt{mpmath.iv} arithmetic, rather than a proof-assistant formalization.
\end{theorem}
We prove the endpoint $L=3/2$; the other windows are its orthogonal compressions in the original $L^2((0,L),\dd x)$ formulation. The zero part of the operator does not affect the norm. Put $V=2\log(3/2)$ and retain the signed kernel $K$ above.

\subsection*{Fourier dilation and an exact finite Schur complement}
Extend $K|_{(0,V)}$ by zero to the circle of length $2V$. Its Fourier multipliers are
\[
 \lambda_k=\int_0^V K(v)e^{-i\omega_kv}\,\dd v,
 \qquad \omega_k=\frac{\pi k}{V}.
\]
There is no normalization divisor in $\lambda_k$. The periodic Hankel operator is the Fourier multiplier followed by reflection. Restriction to $[-V/2,V/2]$ recovers the desired operator exactly: negative sums lie in the zero half of the circle. In the normalized real basis consisting of the constant, cosines and sines, the multiplier-reflection operator has the blocks
\[
 D_0=\lambda_0,\qquad
 D_k=\begin{pmatrix}\Rea\lambda_k&-\Im\lambda_k\\-\Im\lambda_k&-\Rea\lambda_k\end{pmatrix},
 \quad \|D_k\|=|\lambda_k|.
\]
Keep $|k|\leq16$, so the finite matrix $D$ has dimension 33. Let $W$ restrict these orthonormal circle modes to the active interval and put $G=W^*W$. We prove below that all omitted blocks have norm less than $\delta=17/20$. Since restriction is a contraction, for each sign $s=\pm1$,
\begin{equation}\label{eq:majorant}
 s\mathsf H\preceq \delta I+W(sD-\delta I)W^*.
\end{equation}
This retains the low-frequency signed form exactly and charges the tail only for the energy in omitted modes.

Set $\alpha=1-\delta=3/20$, $B_s=sD-\delta I$, and
\begin{equation}\label{eq:inertia}
 S_s=B_s^{-1}-G/\alpha.
\end{equation}
Every $B_s$ is certified invertible. The block matrix with entries $\alpha I,W,W^*,B_s^{-1}$ has these two Schur complements, giving
\[
 n_-(\alpha I-WB_sW^*)=n_-(S_s)-n_-(B_s).
\]
The nullities also agree. Thus equal negative counts and invertibility of $S_s$ prove strict positivity of the right side of $I-s\mathsf H\succeq\alpha I-WB_sW^*$. No bound on an uncomputed cross block is being assumed.

The Gram matrix is explicit. In the complex basis, $G_{kk}=1/2$ and
\[
 G_{kj}=\frac{\sin(\pi(k-j)/2)}{\pi(k-j)}\quad(k\ne j).
\]
The real cosine and sine Gram blocks are obtained by taking their sums and differences. All sine numerators are exactly $0$, $1$, or $-1$.

\subsection*{Analytic control of every omitted mode}
Write $f(t)=e^{-t/2}g_a(e^{-t})$ and $d=(2\pi)^a/\Gamma(a)$, here with $a=1/4$. For $t>0$,
\[
 f(t)=d\,t^{a-1}F(t)-P(t),\quad
 F(t)=e^{-3t/2}\left(\frac{t}{\sinh t}\right)^{1-a},\quad
 P(t)=\frac{2a\pi^a}{\Gamma(a)}e^{(1/2-a)t}
 B_{1-e^{-2t}}(a,3/2-a).
\]
The function $P$ is increasing and $P(0)=0$. For $0<t\leq T$,
\[
 -F'(t)\leq L_T,\quad 1-F(t)\leq L_Tt,
 \qquad L_T=\frac32+\frac{(1-a)T}{3}.
\]
Indeed, $0<F\leq1$ and $\coth t-1/t\leq t/3$; the latter follows, for example, by differentiating the power series for $\sinh t$ and comparing coefficients. Consequently the absolutely continuous remainder $r(t)=f(t)-dt^{a-1}$ satisfies $r(0)=0$ and
\[
 \operatorname{TV}_{[0,T]}r\leq d(2-a)L_T\frac{T^a}{a}+P(T)=:E_T.
\]
Subtract the full Fourier transform $d\Gamma(a)(i\omega)^{-a}$ of the leading power, then integrate both remainders by parts. Combining their endpoint terms before taking absolute values gives
\begin{equation}\label{eq:tail}
 \left|\int_0^T f(t)e^{-i\omega t}\,\dd t\right|
 \leq (2\pi)^a|\omega|^{-a}
 +\frac{|f(T)|+E_T+dT^{a-1}}{|\omega|}.
\end{equation}
The improper power transform is justified by Abel damping or ordinary oscillatory convergence; its tail derivative is integrable.

The two arithmetic terms have $T_1=V$, $T_2=V-\log2$, and weights $1$ and $\rho=(2^{1/4}-2^{-1/4})/\sqrt2$. Summing \eqref{eq:tail}, including the unit-modulus phase of the second term, proves
\[
 |\lambda_k|\leq A|\omega_k|^{-1/4}+B|\omega_k|^{-1},\qquad
 A<1.973173,\quad B<8.583348.
\]
The executable certificate encloses the unrounded expressions directly. At $|k|=17$ the bound is less than $0.822974956342<17/20$, and it decreases for larger $|k|$. This proves the infinite tail estimate used in \eqref{eq:majorant}.

\subsection*{Validated coefficients and matrix inertia}
No oscillatory numerical quadrature is required for the finite coefficients. Put $C=2\pi^a/\Gamma(a)$, $Q=1-e^{-2T}$,
$p=a-1/2+i\omega$, and $r=(5/2-a+i\omega)/2$. Integration by parts gives the exact identity
\[
 \int_0^T f(t)e^{-i\omega t}\,\dd t
 =C\left[\left(\frac12-\frac ap\right)B_Q(a,r)
       +\frac ap e^{-pT}B_Q(a,3/2-a)\right].
\]
Here $p\ne0$, including at $\omega=0$. For complex $r$, evaluate
\[
 B_Q(a,r)=Q^a\sum_{j\geq0}\frac{c_jQ^j}{a+j},\qquad
 c_0=1,\quad c_j=c_{j-1}\frac{j-r}{j}.
\]
After degree $M=512$, a rigorous tail radius is
\[
 \frac{|c_{M+1}|Q^{a+M+1}}{(a+M+1)(1-\eta)},\qquad
 \eta=Q\max\left(1,\left|1-\frac{r}{M+2}\right|\right)<1.
\]
To see this, $|1-r/n|^2$ is a convex quadratic in $1/n$; its maximum for $n\geq M+2$ is attained at an endpoint. The remaining denominator ratio is at most one. Both real and imaginary tail errors are enclosed by this radius.

At 80 decimal digits, the script builds interval matrices $S_s$ and obtains candidate real congruence matrices $Q_s$ from an ordinary eigensolver. Each binary floating-point entry of $Q_s$ is then embedded exactly in interval arithmetic. The script verifies $\|Q_s^TQ_s-I\|<1$, hence invertibility, and strict signed diagonal dominance of $Q_s^TS_sQ_s$ by interval Gershgorin bounds. Thus the eigensolver is only a source of candidates; none of its asserted eigenvalues is trusted.

The independently reproduced certificate gives
\[
 \begin{array}{c|cc|cc}
 s&n_-(B_s)&n_-(S_s)&\min|\operatorname{spec}S_s|\text{ lower bound}
 &I-s\mathsf H\text{ gap lower bound}\\\hline
 +1&28&28&0.0009903747686385&0.0000222801224437\\
 -1&29&29&0.07828557549575&0.00174098138915
 \end{array}
\]
The smallest singular value of either $B_s$ exceeds $0.01372913754369$. For completeness, if $g$ bounds the spectral distance of $S_s$ from zero, the Woodbury formula and $\|W\|\leq1$ give
\[
 (\alpha I-WB_sW^*)^{-1}
 =\alpha^{-1}I+\alpha^{-2}WS_s^{-1}W^*,\qquad
 \alpha I-WB_sW^*\succeq\frac{\alpha^2g}{1+\alpha g}I.
\]
Both displayed gaps exceed $1/50000$, proving Theorem~\ref{thm:signed}.

\textbf{Scope of the advance.} The first-prime signed window is now certified, despite its absolute-kernel norm being greater than one. The reusable mechanism is an exact Fourier dilation, an analytic tail bound, and a finite verified Schur complement. This resolves one fixed width and a bounded truncation range. Uniform estimates as $a\downarrow0$, arbitrarily large $L$, and the global RH implication remain unproved. Literature novelty has not been assessed.

\section*{Verification and sources}
The accompanying script checks the harmonic dimensions and derivatives exactly with rational symbolic arithmetic, verifies the completion value and cocycle numerically at 70 decimal digits, and checks finite cyclic-root identities. These are reproducibility diagnostics; the proofs above establish the stated theorems. No claim of literature priority is made.

\begin{enumerate}[label={[\arabic*]}]
\item \emph{Split Twistor Bridge V2}, 1 October 2026, pp. 2--9. Existing project report containing the exact homogeneous Fourier family, prime factors, incidence map at the self-dual weight, and constant-mode reduction. This continuation does not replace that report.
\item J. C. Lagarias, \emph{Zero Spacing Distributions for Differenced L-Functions}, Acta Arithmetica 120 (2005), 159--184; corrected preprint, Lemma 2.1. \url{https://arxiv.org/abs/math/0601653}.
\item M. Suzuki, \emph{A canonical system of differential equations arising from the Riemann zeta-function}, revised 2016. \url{https://arxiv.org/abs/1204.1827}.
\item G. \'Olafsson, A. Pasquale and B. Rubin, \emph{Analytic and Group-Theoretic Aspects of the Cosine Transform} (2012). Background on Fourier normalization and generalized principal-series intertwiners. \url{https://arxiv.org/abs/1209.1822}.
\end{enumerate}
\end{document}
